El presente Trabajo de Fin de Grado se propuso alcanzar una serie de objetivos para el desarrollo y la implementación de la aplicación web ``LogArt: Galería de Recuerdos''. Estos objetivos se detallan a continuación:

\section{Objetivo Principal}
\label{sec:objetivoPrincipal}

Diseñar, desarrollar e implementar una aplicación web funcional, denominada LogArt, que permita a los usuarios registrar, organizar y consultar de manera personal y detallada sus recuerdos, notas y experiencias relacionadas con diversas disciplinas culturales como la música, los videojuegos y la literatura.

\section{Objetivos Específicos}
\label{sec:objetivosEspecificos}

Para conseguir el objetivo principal, se establecieron los siguientes objetivos específicos:

\begin{itemize}
    \item \textbf{OE1: Análisis y Definición de Requisitos:}
    Investigar y definir las funcionalidades clave que la aplicación LogArt debía ofrecer para satisfacer la necesidad identificada de almacenar y gestionar recuerdos culturales.
    \item \textbf{OE2: Diseño de la Arquitectura del Software:}
    Diseñar una arquitectura robusta y escalable para la aplicación, seleccionando un stack tecnológico adecuado (MERN: MongoDB, Express.js, React, Node.js) que facilitase el desarrollo y el mantenimiento a futuro.
    \item \textbf{OE3: Desarrollo del Backend:}
    Implementar la lógica de negocio y la API RESTful en el servidor (Node.js con Express.js) para gestionar los datos de los usuarios, las categorías de disciplinas, las obras y las entradas de recuerdos, asegurando la persistencia de los datos en una base de datos NoSQL (MongoDB).
    \item \textbf{OE4: Desarrollo del Frontend:}
    Construir una interfaz de usuario intuitiva y atractiva con React, que permita a los usuarios, entre otras cosas, registrarse, iniciar sesión, navegar por las diferentes disciplinas, añadir nuevas obras, crear entradas de recuerdos detalladas y visualizar sus registros.
    \item \textbf{OE5: Implementación de Funcionalidades Clave (entre otras):}
        \begin{itemize}
            \item Gestión de usuarios (registro, inicio de sesión, perfil).
            \item Creación, visualización, edición y eliminación (CRUD) de Obras para diferentes tipos de disciplinas (libros, videojuegos, canciones), permitiendo la subida de imágenes.
            \item Creación, visualización, edición y eliminación (CRUD) de entradas de recuerdos para las diferentes obras.
        \end{itemize}
    \item \textbf{OE6: Pruebas y Validación:}
    Realizar pruebas unitarias, de integración y End-to-End (E2E) para asegurar la calidad, el correcto funcionamiento de la aplicación y la ausencia de errores críticos.
    \item \textbf{OE7: Empaquetado y Despliegue:}
    Empaquetar la aplicación utilizando Docker y Docker Compose para simplificar su configuración, despliegue y portabilidad, y configurar un flujo de Integración Continua y Despliegue Continuo (CI/CD) utilizando GitHub Actions.
    \item \textbf{OE8: Implementación de Funcionalidades Avanzadas (entre otras):}
        \begin{itemize}
            \item Recuperación de contraseña a través de email, y verificación de cuenta del mismo modo.
            \item Panel de administrador para visualizar diferentes tipos de gráficos y datos sobre el desempeño de la aplicación.
            \item Desarrollo de los test para las funcionalidades avanzadas.     
        \end{itemize}
    \item \textbf{OE9: Elaboración de la Documentación:}
    Redactar la presente memoria de Trabajo Fin de Grado, documentando todo el proceso de desarrollo, las decisiones tomadas, la arquitectura del sistema y los resultados obtenidos.
\end{itemize}

El cumplimiento de estos objetivos permitiría obtener una primera versión completa y funcional de LogArt, sentando las bases para posibles futuras mejoras y expansiones.